\documentclass[1p,times]{elsarticle}
\usepackage[utf8]{inputenc}
\usepackage[T1]{fontenc}
\usepackage{graphicx}
\usepackage{amsmath,amssymb}
\usepackage{booktabs}
\usepackage{lineno}
\usepackage[hidelinks]{hyperref}
\graphicspath{{.}{fig/}{../figure/set_v2/}}
\biboptions{numbers,sort&compress}
\journal{International Journal of Disaster Risk Reduction}
\linenumbers

\usepackage{url}
\begin{document}

\begin{frontmatter}

\title{How much of hazard news is news? Decomposing open-source coverage of chemical accidents and CBRN threats for disaster risk reduction}

\author[cbrn]{Miae Kim}
\author[knpa]{Janghyeon Cha}
\author[cbrn]{Myeongsik Shin}
\author[cbrn]{Gyungjun Choi}
\author[cbrn]{Ku Kang\corref{cor}}
\ead{bisu9082@gmail.com}
\cortext[cor]{Corresponding author.}
\address[cbrn]{Chemical, Biological, Radiological and Nuclear (CBRN) Defense Research Institute, Seoul, Republic of Korea}
\address[knpa]{Korean National Police Agency, Seoul, Republic of Korea}

\begin{abstract}
Planners read news volume as a proxy for hazard frequency, especially for chemical, biological, radiological and nuclear (CBRN) threats, which have no registry. Is an article about a hazard evidence that anything happened? We collected 109,610 Korean articles (54 outlets, 2021--2026) for 30 queries in three blocks: industrial chemical accidents, CBRN security threats and conventional terrorism. A language model coded a stratified gold set of 2,867 articles into five components from new event to off-topic, a second model coded 600 of them for reliability, and two authors then coded a 400-article subset independently ($\kappa = 0.85$). The models matched them for conventional terrorism and diverged severely for CBRN threats ($\kappa = 0.24$), counting topically adjacent occurrences the codebook excludes. On the human standard, 219 agreed articles in this block, 5.0\% of CBRN security articles report an event (95\% CI 2.5--7.9; 4.4--6.0 across treatments of coder disagreement) against 38\% and 34\% in the controls: 22.5 articles per unique CBRN event (95\% CI 14.2--44.0), a lower bound by design. Excluding off-topic retrievals leaves 13.9\% against 52.1\% and 59.6\%, so the contrast is not an artefact of query precision; its significance tests use model labels, which overstate the CBRN share but not the contrast. Reference bias, the share of articles citing a past anchor event, is 89\% for Novichok and 33\% for hydrogen fluoride against 1--2\% for same-sector controls, and a rule naming no anchor reproduces the gap because the controls have no canonical occurrence to cite. Article counts for six chemicals bear no detectable relation to registry counts (log--log slope $-0.07$, CI $-0.89$ to $0.76$). The decomposition converts volume into a bounded event estimate and separates inherited salience from current incidence; we set out the procedure an agency can run on its own archive.
\end{abstract}

\begin{keyword}
event-based surveillance \sep CBRN \sep chemical accidents \sep social amplification of risk \sep news media \sep reference bias \sep disaster risk reduction
\end{keyword}

\end{frontmatter}

%% ====== HIGHLIGHTS (<=85 chars each; separate file highlights.tex) ======
%% - News volume for 30 hazard queries is decomposed into five labelled components
%% - 11% of CBRN security coverage reports an event, and 2% reports a domestic event
%% - Reference bias reaches 33-89% for rare hazards and 1-2% for matched controls
%% - Article counts for six chemicals are unrelated to official incident counts
%% - A lightweight classifier recovers accident and terrorism events; CBRN needs coded samples

%% ------------------------------------------------------------------
\section{Introduction}
\label{sec:intro}

Agencies that plan for chemical, biological, radiological and nuclear (CBRN) incidents rarely have an incident registry to size the threat. For industrial chemical accidents a registry exists in most countries; for anthrax letters, dirty bombs or nerve-agent attacks it does not, because the events are too rare. Detection technology for such agents has advanced quickly, but sizing the threat that the technology is meant to meet is still a question answered from open sources. In the authors' own institutional setting, threat briefings count articles, trend lines of coverage are read as trend lines of risk, and detection and protection assets are argued for on the strength of how much a hazard is talked about. We cannot document how widely that practice holds elsewhere, because such briefings are rarely published; what is documented is the wider family it belongs to, in which media reports are treated as an indicator of hazard activity for surveillance and early warning \cite{hartley2010,barboza2014,freifeld2008,brownstein2009,abbood2020}. The question this paper asks is prior to how common the practice is: whether the count, used that way, carries the information it is read as carrying. The practice is understandable, since media salience shapes both public and expert judgement of what is frequent \cite{mccombs1972,tversky1973,combs1979,slovic1987}, and the Republic of Korea's public-health service has shown that open-source event surveillance can work \cite{kim2025ebs}. It rests, however, on an assumption that is seldom tested: that an article about a hazard is evidence that something happened.

This study tests that assumption. We treat the news volume $V$ observed for a hazard query as a sum of components that have different meanings for a planner. Some articles report a new event. Some repeat an event already reported. Some cite a past event (the 2012 Gumi hydrogen fluoride leak \cite{kim2012gumi,koo2017gumi}, the 1995 Tokyo subway sarin attack \cite{okumura1996}, the 2018 Salisbury and 2020 Navalny Novichok poisonings \cite{chai2018}) without any new event having occurred. Some describe drills, inspections, legislation and budgets. Some are on-topic but report nothing, and some are retrieval errors. Only the first component says anything about incidence. The rest inflate $V$ in ways that differ by hazard, so two hazards with equal coverage can have very different event counts.

Several literatures bear on this problem. Sociologists who use newspapers to count protest events distinguished selection bias (which events get covered) from description bias (how they are described) four decades ago \cite{earl2004,strawn2008,ortiz2005,zhang2023}, and the builders of conflict and terrorism event datasets have documented how reporting bias shapes what enters a database \cite{weidmann2016,baum2015,raleigh2010,lafree2007,wang2016}; in these frameworks the unit is the event, and an article that cites a past event without reporting a new one is not a category. Communication research has described how key events reshape subsequent coverage of a hazard class \cite{kepplinger1995,brosius1995}, how past events serve as templates through which new ones are reported \cite{kitzinger2000}, how journalists draw on collective memory \cite{edy1999}, how availability cascades link salience to regulation \cite{kuran1999}, and how news waves become self-reinforcing \cite{vasterman2005,wien2009}; these studies establish the mechanism but do not give a planner a ratio to compute for a given hazard. The social amplification of risk framework \cite{kasperson1988,pidgeon2003,kasperson2022} names the media as an amplification station, and its quantitative applications relate coverage volume to perceived risk and societal impact \cite{renn1992,burns1993}; coverage is known to move public risk perception \cite{wahlberg2000,nellis2012} and even expert risk perception \cite{won2025}. Event-based surveillance in public health uses news as a signal for outbreaks, poisonings and hazardous exposures \cite{freifeld2008,brownstein2009,hartley2010,barboza2014,abbood2020,mcknight2024,cook2021,rodde2025}, and disaster-risk researchers mine news and social media for damage and response assessment \cite{olteanu2015,imran2015,murthy2025,vescovo2026,erokhin2024}. Operational epidemic intelligence is not the target of this paper, and it is worth being exact about why. Systems of the EIOS family separate a signal from a verified event by procedure: articles are triaged, a signal is verified against another source, and only then does it become an event for risk assessment \cite{williams2025eios,rodde2025}. Their published evaluations report the yield of that pipeline, which is the quantity this paper computes for a hazard's coverage as a whole. What we address is the step before any such pipeline, and the practice that has no pipeline at all: reading the volume of retrieved articles, or its trend, as a measure of how much is happening. The history of search-volume epidemiology \cite{ginsberg2009,lazer2014} shows what happens when a proxy is trusted without a ground-truth comparison.

We make four contributions. First, we state an identity that decomposes $V$ into five labelled components and define an amplification factor $A$, the number of articles per unique event, that converts volume into an event estimate. Second, we define reference bias $\rho$ as the share of articles that cite a past anchor event and measure it against same-block controls scored on the same anchor. Third, we compare news-derived counts with the national chemical-accident registry and show that article counts are unrelated to incident counts. Fourth, we show that a character $n$-gram classifier recovers the event component for accident and terrorism queries, including queries it was not trained on, but not for CBRN queries, where events are too few and too varied, and that the archive's own incident taxonomy does not help, because it encodes what an article is about rather than whether anything happened.

The setting is Korean national news from 54 outlets over five years. The three-block design, with industrial accidents that have a registry, CBRN security threats that do not, and conventional terrorism as a second control, lets us ask whether CBRN coverage is unusual, not only whether it is inflated. Which attacks and accidents receive coverage depends on the perpetrator, the casualty count and the hazard's novelty \cite{kearns2019,eisensee2007}; the design holds outlet set and window fixed so that these selection effects are compared across blocks rather than assumed away.

%% ------------------------------------------------------------------
\section{Decomposition of news volume}
\label{sec:framework}

Let $V$ be the number of articles retrieved for a query over a fixed window and outlet set. We partition every article into one of five mutually exclusive components (Fig.~\ref{fig:framework}a):
\begin{equation}
V = \mathrm{EVENT} + \mathrm{REF} + \mathrm{INST} + \mathrm{OTHER} + \mathrm{IRRELEVANT}.
\label{eq:identity}
\end{equation}
EVENT articles have a newly occurring event as their main subject, whether first report or follow-up. REF articles have a past event as their main subject with no new occurrence. INST articles report institutional activity: drills, inspections, legislation, budgets, agreements and official statistics. OTHER articles are on topic but report neither an event nor an institutional act. IRRELEVANT articles were retrieved by the query but concern something else.

\begin{figure*}[!t]
\centering
\includegraphics[width=\textwidth]{Fig1_framework_data.pdf}
\caption{Framework and corpus. (a) Decomposition of observed news volume $V$ into five components, with EVENT shown as NEW + FOLLOW; $N_{\mathrm{events}} = \mathrm{EVENT}/m$ and $A = V/N_{\mathrm{events}} = m/p_E$. Reference bias $\rho$ is measured separately as the share of articles citing an anchor event. (b) Articles per query for 30 queries in three blocks (BigKinds, 54 outlets, 9 Sep 2021 to 8 Sep 2026; log scale). (c) Monthly rate of ``drone attack'' articles by year; 2021 and 2026 are partial years.}
\label{fig:framework}
\end{figure*}

EVENT articles are grouped into unique events by date, place and object. If $m$ is the mean number of EVENT articles per unique event, the event count is $N_{\mathrm{events}} = \mathrm{EVENT}/m$ and the amplification factor is
\begin{equation}
A = \frac{V}{N_{\mathrm{events}}} = \frac{m}{p_E}, \qquad p_E = \frac{\mathrm{EVENT}}{V}.
\label{eq:A}
\end{equation}
$A$ is the factor by which raw volume overstates incidence; it is an arithmetic ratio and is not the amplification of the social amplification framework, which concerns perceived risk and secondary impacts. Its two parts behave differently under sampling: $p_E$ is a proportion that a random sample estimates without bias, whereas $m$ can only be observed for events whose reports fall into the sample together, so $m$ estimated from a small sample of a large query is biased towards one. We therefore treat $p_E$ as the primary quantity for comparing hazards, report $m$ separately, and treat $A$ from a sample as a lower bound on the corpus-internal ratio. That direction holds only for the observed-event denominator; differential capture pushes the other way and is not bounded by these data (Section~\ref{sec:capture}).

Reference bias is measured separately from the labels. A REF label requires that the past event be the article's main subject, which excludes the common case in which a new event is reported and a past event is invoked in passing. We therefore define
\begin{equation}
\rho = \frac{\#\{\text{articles mentioning a past anchor event}\}}{V}
\label{eq:rho}
\end{equation}
for hazards with an identifiable anchor, and score control hazards from the same block against the same anchor so that outlet, period and sector are held constant. Because an anchor lexicon can also match new developments of the anchor event itself, we report $\rho$ by anchor and by year and cross-tabulate anchor mentions with the EVENT label in the gold set.

%% ------------------------------------------------------------------
\section{Materials and methods}
\label{sec:methods}

\subsection{Corpus}
Articles were retrieved from BigKinds, the Korea Press Foundation's news archive, which indexes 54 national and regional outlets. Thirty queries were fixed before collection, ten per block, from the substance list of the national chemical-accident registry (industrial), the threat lexicon used in the authors' institutional threat assessments (security) and the incident categories of the police counter-terrorism classification (conventional); no query was added or dropped after collection (Fig.~\ref{fig:framework}b): industrial chemical accidents (e.g.\ hydrogen fluoride leak, ammonia leak, chemical plant explosion), CBRN security threats (e.g.\ chemical weapon use, anthrax terror, dirty bomb, Novichok, DPRK nuclear test) and conventional terrorism (e.g.\ explosive terror, drone attack, mass shooting, bomb threat). One boundary query, ``reactor radiation exposure'', was placed in the conventional block by design as a hazard with a domestic institutional base; its assignment is tested in Section~\ref{sec:sens}. The window was 9 September 2021 to 8 September 2026, the maximum five-year span the archive allows with a search term, applied identically to all queries; all downloads were made on 8 September 2026. Queries exceeding the archive's 20,000-record export limit were retrieved by calendar year and de-duplicated on the archive's article identifier. Because BigKinds treats a space as a logical AND, the two-token query ``nerve agent Novichok'' returned zero records and was replaced by ``Novichok''. The corpus contains 109,610 articles: 4,633 industrial, 29,688 security and 75,289 conventional-terrorism. Each record carries title, the first 200 characters of body text, outlet, date, the archive's named locations and its topic and incident taxonomies. The 200-character limit is the archive's copyright policy; all coding and modelling below use title plus this excerpt only.

Because the archive names every export file identically, file--query assignment was verified by content rather than by download order: for each file, the share of records containing all query tokens in title, keywords or excerpt was computed against every query, and the file was assigned to the query with the highest share. Assignments were unambiguous (typically 1.00 against 0.00) and the procedure caught nine files that a download-order assignment had mislabelled; the nine files and both assignments are listed in the supplementary manifest.

\subsection{Gold set and coding}
\label{sec:coding}
A gold set of 2,867 articles was drawn by stratified random sampling with query by year as strata (178 non-empty strata). Each query was allotted 100 articles apportioned to years in proportion to stratum size with at least three per year; three queries with fewer than 100 articles (sarin, 41; white-powder report, 59; hazardous chemical release, 67) were taken in full. The sampling seed is fixed (20260908) and the sample was shuffled before coding. The sampling fraction $n/V$ ranges from 0.2\% (``drone attack'') to 100\%; eight queries have fractions below 5\% and sixteen above 20\%.

A codebook defines the five components with a decision order (IRRELEVANT, EVENT, REF, INST, OTHER) and rules for the boundary cases: an article is judged by its main subject, so a new event that invokes a past one is EVENT; follow-up coverage of an event whose handling is still in progress (rescue, investigation, prosecution) is EVENT, whereas retrospective coverage after the process has closed (anniversaries, memoirs, policy reviews) is REF; official incident statistics count as INST. The codebook was revised once after a 100-article pilot coded by the two model coders described below ($\kappa = 0.772$), in which all disagreements fell between non-event categories. For EVENT articles coders also recorded an estimated event date, place and object.

Two large language models were used as independent coders under the same codebook, following the emerging practice of model-based annotation in computational social science \cite{gilardi2023,ziems2024} and its documented caveats about reproducibility with proprietary models \cite{ollion2024} and the need for validation against human labels \cite{pangakis2023}. Coder A (Claude Opus, Anthropic) coded all 2,867 articles in ten batches of about 290 articles, each article as a separate item with no access to other articles' labels; coder B (Claude Sonnet, Anthropic) coded a random subset of 600 (232 industrial, 181 security, 187 conventional) without access to A's labels. The codebook instructs the coder to judge from title and excerpt only. Both were run at temperature 0 on 8 September 2026 with the codebook as the complete instruction. The codebook was at version 2, revised after a 100-article pilot, and is deposited with the labels. We departed from the deposited protocol in three respects. It specifies double coding of all 2,867 articles, whereas coder B coded a random 600 for reliability estimation only, so the reliability gate is met on those 600 rather than on the full gold set. It specifies a random 400-article human subset, whereas the subset drawn was disproportionately stratified, taking every article coder A called an event in the security block; the design is described below and the estimates are post-stratified accordingly. And it specifies adjudication of disagreements, which was performed for neither the 94 model--model disagreements nor the 40 human disagreements; for the human ones the effect is reported as a range instead, and for the model ones the labels are used as coder A gave them. The protocol's reliability gate ($\kappa \ge 0.70$) was met without revision to a third version. The wrapper text that presented each article to the model and the model identifiers returned by the interface were not retained, which limits exact reproduction of the coding runs. Version 2 contains the rule that decides most of the model--human divergence reported in Section~\ref{sec:human}: an article is off-topic when its subject is a different agent or method from the one the query names, so that a nuclear-weapons report retrieved by a chemical-weapons query is IRRELEVANT rather than an event. Both model coders and both human coders worked from that same text. The two coders come from the same vendor, so their agreement bounds reliability from above and is not evidence of validity. Agreement is reported as Cohen's $\kappa$ \cite{cohen1960} with an asymptotic 95\% confidence interval and interpreted on the conventional scale \cite{landis1977,mchugh2012}. A 400-article subset was coded independently by two of the authors, denoted H1 and H2, blind to the model labels and to each other, to establish validity against human judgement. The security block was sampled disproportionately, taking all 95 articles coder A had labelled EVENT together with 150 drawn at random from the remaining 805, so that the block carrying the study's main claim contributes a census of its putative events; the industrial and conventional blocks were sampled at random within the block (78 and 77 articles). Estimates on the subset are post-stratified onto the six cells formed by block and coder A's EVENT-versus-rest decision, whose sizes in the gold set are known, and intervals come from a within-cell bootstrap of 10,000 resamples. The reference standard is the set of articles the two coders agreed on, which is a non-random deletion of the 40 they did not: the post-stratification weights reflect the probability of being sampled, not the probability of being agreed on, and agreement is itself related to the boundary the estimate turns on. We therefore report the estimate across treatments of the disagreements rather than the consensus value alone, and treat the width of that range, not the bootstrap interval, as the statement of what the figure is known to within. Coder A's sensitivity and false-positive rate against this standard are also reported, and are used to carry the correction to the full corpus by the Rogan--Gladen estimator \cite{rogan1978}; on the gold set itself that estimator is algebraically identical to the post-stratified estimate under the same weights and is therefore not reported as a separate check. The result is reported in Section~\ref{sec:human}. Unique events were formed from EVENT articles with matching place and object and dates within three days; this entity-and-window rule was chosen over embedding-based article matching \cite{sen2019} so that the pipeline runs without a GPU, and its sensitivity to the window and to dropping place or object is reported.

\subsection{Statistics}
\label{sec:stats}
Block-level component shares are reported from the gold set with 95\% confidence intervals from a stratified bootstrap (10,000 resamples within query-by-year strata) \cite{efron1979}. The amplification factor is reported with $m$ held at its sample value and only $p_E$ resampled, and is labelled a lower bound for the reason given in Section~\ref{sec:framework}.

The main test of block differences uses the query as the unit, since articles within a query are not independent, and uses the event share $p_E$ rather than $A$, since $m$ depends on the sampling fraction. The 30 queries were drawn from three external lists, not sampled at random from a population of queries, so every block comparison is an inference about these 30 and their blocks, not about CBRN queries in general; block-level shares, by contrast, are design-based estimates for the corpus these queries retrieved. A Kruskal--Wallis test with $\eta^2$ and $\epsilon^2$ and pairwise Mann--Whitney tests (asymptotic, with Holm-adjusted $p$ and rank-biserial $r$ \cite{kerby2014}) were specified in advance for ten queries per block; exact $p$-values are smaller and are given in the supplement. Article-level differences are described with Cram\'er's $V$ and Cohen's $h$ without a $p$-value, because articles are clustered within queries.

\subsection{Domestic and foreign events}
For every EVENT article the coded place was classified as domestic (Korean administrative unit or nationwide) or foreign (foreign country or city); the 47 articles with no coded place were left unknown, and 56 articles whose coded place matched neither lexicon (mostly foreign regions and organisations) were classified as foreign by the same rule applied to the title and excerpt. Places in the DPRK were classified as foreign, and a foreign country or city named in the place field took precedence over a Korean-style administrative suffix. The 18 domestic security-block EVENT articles are listed in the supplement. This distinction matters because a planner's question concerns domestic incidence, whereas Korean coverage of terrorism and CBRN threats is largely foreign.

\subsection{Anchor events and official registry}
For five hazards with a widely recognised anchor event, $\rho$ was computed as the share of articles whose title, keywords or excerpt mention the anchor. Lexicons, with spelling variants, were: Novichok (Skripal, Salisbury; Navalny), sarin (Aum Shinrikyo, Tokyo; Ghouta), radiation leak (Fukushima; Chernobyl), anthrax terror (United States, mail letters \cite{jernigan2002}) and hydrogen fluoride leak (Gumi). Bare year tokens (2012, 1995, 2001, 2018, 2020) were scored separately and excluded from the headline $\rho$, because they also match unrelated dates. Three industrial queries (chlorine, sulfuric acid and nitric acid leaks) were scored against the Gumi anchor as same-block controls. Nuclear accidents were included because their coverage is known to be dominated by a small number of canonical events \cite{friedman2011}. The archive supplies for every article a keyword field extracted automatically from the full text, which the coders did not see. We therefore report $\rho$ on two bases: title and excerpt alone, and title, excerpt and keywords. The gold-set cross-tabulation with the EVENT label is computed on the second basis by joining the keyword field to the sampled articles.

The Ministry of Environment's chemical-accident registry (CARIS), maintained by the National Institute of Chemical Safety, was obtained from the national open-data portal (1,026 incidents, January 2014 to April 2025). Comparison used the overlap window 9 September 2021 to 28 April 2025 (366 incidents). The registry underpins Korea's off-site emergency planning and has been analysed for response and disclosure purposes \cite{nam2025,yoon2025}, but not, to our knowledge, against news coverage. Six substance queries were matched to registry incidents by substring search on the registry's three substance fields, which are free text rather than a controlled vocabulary; the rule for each substance is deposited. The six matched sets contain 142 distinct incidents and 151 substance assignments, nine incidents being mixed acids that carry two of the six substances at once. Article counts in the same window were regressed on incident counts in log--log form; the practice under test corresponds to a slope of 1, so that hypothesis is tested as well as the null of no relation. The framework's own output was also compared with the registry: for each substance the estimated number of domestic events is $\hat N = V p_E^{\mathrm{corpus}} d / m$, where $p_E^{\mathrm{corpus}}$ is the post-stratified event share for the query, $d$ the domestic share of its events and $m$ its sample duplication factor. We regress articles, not articles per incident, on incidents: a regression of a ratio on its own denominator has a slope of $-1$ under the null of no relation \cite{pearson1897}, so the ratio form cannot be used as evidence.

\subsection{Classifier}
\label{sec:clf-methods}
The proposed model is a character $n$-gram (2--4) TF--IDF representation (minimum document frequency 2, sublinear term frequency, at most 120,000 features) of title plus excerpt with a class-weighted multinomial logistic regression, a linear text-categorisation design chosen for transparency and small footprint \cite{sebastiani2002}. The vectoriser is fitted inside each training fold; $C = 4$ was fixed before the experiment and not tuned. Baselines are majority class, uniform random, and the keyword rule that underlies most volume-based practice: if all query tokens occur in title or excerpt the article is EVENT, otherwise OTHER. Ablations remove the excerpt, replace character with word $n$-grams, and add a one-hot encoding of the archive's incident taxonomy. Evaluation uses ten random seeds (0--9) by five-fold cross-validation stratified on the label (50 folds); the primary metric is macro-$F_1$ with $F_1$ for the EVENT class reported alongside, and per-block scores are reported for the proposed model. Because folds that reuse the same 2,867 articles are not independent, model differences are tested with the Nadeau--Bengio corrected resampled $t$-test \cite{nadeau2003,dietterich1998} ($J = 50$ folds, test-to-train ratio $0.25$) and reported with seed-level 95\% confidence intervals ($n = 10$), which reflect split variance only. Because the cross-validation folds share all 30 queries between training and test, so that a character model can learn query vocabulary and with it the query's base rate, the model was also evaluated leave-one-query-out, training on 29 queries and predicting the thirtieth. Three sanity checks were specified before the experiment, the third for macro-$F_1$ only: macro-$F_1$ must rise monotonically with training fraction (25, 50, 75, 100\%); random-baseline accuracy must equal $1/K$; and per-block macro-$F_1$ must not diverge by more than a factor of 1.5.

\subsection{Full-corpus estimates}
Because the three blocks' queries were written from three different sources, they do not retrieve with equal precision, and a share of retrieved articles is partly a statement about the query. We therefore report the event share twice: unconditionally, $p_E = \mathrm{EVENT}/V$, and conditional on topical relevance, $p_E^{\mathrm{rel}} = \mathrm{EVENT}/(V - \mathrm{IRRELEVANT})$, which removes the retrieval difference and asks what share of on-topic coverage reports an event. The second is the quantity that isolates the hazard from the query.

Two quantities were estimated for the full corpus. The event share was estimated by post-stratification, $p_E^{\mathrm{corpus}} = \sum_s V_s \hat p_s / \sum_s V_s$ over the 178 query-by-year strata, which uses the gold labels and the known stratum sizes and needs no classifier; its interval comes from a binomial bootstrap within strata, and a second interval from resampling whole queries within the block, which captures between-query variance and is the more conservative of the two. The duplication factor $m$ was estimated by applying the classifier, trained on the whole gold set, to the corpus and grouping predicted EVENT articles within a query when they fell within three days and their title-plus-excerpt TF--IDF cosine similarity was at least 0.25. The similarity threshold was chosen from $\{0.20, 0.25, 0.35\}$ with and without the location constraint by maximising the adjusted Rand index against the coders' clusters on the gold-set EVENT articles, so it is tuned on the same labels against which it is then compared and its agreement is an upper bound. On that comparison it overshoots the aggregate $m$ by a tenth (1.46 against 1.32, events counted within queries) and agrees with the coders on only a minority of individual article pairs (pairwise $F_1 = 0.36$); across the six settings the text-rule $m$ ranges from 1.24 to 1.59. The corpus $m$ is therefore reported as indicative, with that range, and is not used where the classifier's event predictions are themselves unreliable. The post-stratified corpus shares were additionally corrected for coder A's measured misclassification, using the sensitivity and false-positive rate estimated against the human reference standard of Section~\ref{sec:human}. Uncertainty in the observed share is propagated by drawing it from a normal approximation to its published interval and uncertainty in the rates by resampling the human subset within cells; the two are treated as independent, which is approximate because both use gold-set articles. The correction is informative only where the observed share stands clear of the false-positive rate, and this condition is stated with each corrected value. The classifier's own cross-validated sensitivity and false-positive rates are not used for this purpose.

%% ------------------------------------------------------------------
\section{Results}
\label{sec:results}

\subsection{Composition of coverage}
\label{sec:composition}
Across the gold set, 32.2\% of articles were EVENT (95\% CI 30.7--33.6), 2.9\% REF, 27.8\% INST, 20.9\% OTHER and 16.2\% IRRELEVANT (Table~\ref{tab:components}, Fig.~\ref{fig:decomp}a). The blocks differ sharply. In the industrial block 43.2\% (40.6--45.8) of articles reported an event and 34.5\% reported institutional activity. In the conventional-terrorism block 40.9\% (38.3--43.5) reported an event. In the CBRN security block 10.6\% (8.9--12.3) reported an event; 35.3\% were on-topic non-event material, 25.9\% institutional and 25.6\% off-topic retrievals. The security figure is the one the human validation of Section~\ref{sec:human} revises: against two human coders the security event share is 5.0\% (2.5--7.9), because the model counted occurrences that the codebook excludes as retrievals of the wrong kind. Every value in this section, including the tests that follow, is computed on model labels; the human-anchored event shares are given in Section~\ref{sec:human} and are the ones carried into the discussion. The tests are re-run on the human labels in Section~\ref{sec:human}, where the subset's size per query and what it can and cannot resolve are set out.

The blocks also differ in how cleanly their queries retrieve, and part of the raw contrast is that difference rather than a difference in event rates: 25.6\% of security articles are off-topic against 9.2\% industrial and 14.5\% conventional. Conditioning on topical relevance removes it. Among on-topic articles the event share is 47.6\% in the industrial block, 47.8\% in the conventional block and 14.2\% in the security block, and the query-level test on this quantity gives $H(2) = 9.29$, $p = 0.0096$, $\eta^2 = 0.27$, with security below industrial (Holm-adjusted $p = 0.022$, $r = 0.60$) and below conventional ($p = 0.028$, $r = 0.55$) and the two controls again indistinguishable ($p = 0.62$). A CBRN query that has already returned an on-topic article is still three times less likely to have returned an event than either control; Section~\ref{sec:human} repeats this on human labels. The article-level contrast between the security block and either control is large (descriptively, since articles cluster within queries: Cram\'er's $V = 0.28$ over the five components and 0.31 for EVENT versus the rest; security versus industrial, Cohen's $h = 0.77$), while the two controls are alike ($h = 0.05$).

\begin{table}[!htb]
\centering\small
\caption{Component decomposition of news volume in the gold set ($n = 2{,}867$). CI, 95\% stratified bootstrap interval for $p_E$; $m$, EVENT articles per unique event in the sample (biased towards 1 for queries with small sampling fractions); $A = m/p_E$, a lower bound, with events merged across queries within a block (within-query clustering gives $m = 1.63$, 1.10, 1.15 and 1.32), with CI obtained from the CI of $p_E$ with $m$ fixed. IRREL., irrelevant.}
\label{tab:components}
\setlength{\tabcolsep}{3.5pt}
\begin{tabular*}{\textwidth}{@{\extracolsep{\fill}}lrlrrrrrl@{}}
\toprule
Block & $n$ & EVENT \% [CI] & REF & INST & OTHER & IRREL. & $m$ & $A$ [CI] \\
 & & & (\%) & (\%) & (\%) & (\%) & & \\
\midrule
Industrial & 967 & 43.2 [40.6, 45.8] & 2.2 & 34.5 & 10.9 & 9.2 & 1.75 & 4.05 [3.82, 4.30] \\
Security (CBRN) & 900 & 10.6 [8.9, 12.3] & 2.7 & 25.9 & 35.3 & 25.6 & 1.12 & 10.59 [9.06, 12.57] \\
Conventional terror & 1,000 & 40.9 [38.3, 43.5] & 3.8 & 23.1 & 17.7 & 14.5 & 1.19 & 2.90 [2.73, 3.10] \\
All & 2,867 & 32.2 [30.7, 33.6] & 2.9 & 27.8 & 20.9 & 16.2 & 1.38 & 4.29 [4.11, 4.49] \\
\bottomrule
\end{tabular*}
\end{table}

The pattern is visible query by query. For ``bioterrorism'' one article in 100 reported an event and 74 reported drills or preparedness activity; for ``anthrax terror'' seven in 100 reported an event; for ``DPRK nuclear test'', with no test in the window, the seven EVENT articles report missile launches and test-site activity that the codebook treats as events of the query's object. For ``bomb threat'', by contrast, 69 in 100 reported an event.

At the query level (Fig.~\ref{fig:decomp}c) the event share separates the blocks (Kruskal--Wallis $H(2) = 10.06$, $p = 0.0065$, $\eta^2 = 0.30$, $\epsilon^2 = 0.35$; medians 44\%, 7\% and 50\%). With ten queries per block the effect size is imprecise: resampling queries gives a 95\% interval of 0.09--0.60 for $\eta^2$, and under 10,000 permutations of the block labels $\eta^2$ averages $-0.001$ with a 95th percentile of 0.14, so the observed value is above what this design produces by chance but should be read as a rank separation rather than a variance share. The security block lies below the industrial block (Mann--Whitney $U = 12$, Holm-adjusted $p = 0.014$, $r = 0.76$) and below the conventional block ($U = 16$, $p = 0.022$, $r = 0.68$; $U$ is given for the first-named block), whereas the industrial and conventional blocks are indistinguishable at this sample size ($U = 48.5$, $p = 0.94$, $r = 0.03$).

\begin{figure*}[!t]
\centering
\includegraphics[width=\textwidth]{Fig2_decomposition_amplification.pdf}
\caption{Component decomposition and amplification (gold set, $n = 2{,}867$). (a) Share of the five components by block. (b) Amplification factor $A$ with 95\% CI obtained from the stratified-bootstrap CI of $p_E$ with $m$ held at its sample value; $A$ from a sample is a lower bound. (c) Query-level EVENT share $p_E$ ($n = 10$ per block) with medians; Kruskal--Wallis and pairwise Mann--Whitney tests with Holm-adjusted $p$ and rank-biserial $r$. (d) Sample-based $A$ for each query, $A = n/\text{unique events}$; parentheses give EVENT articles / sample size.}
\label{fig:decomp}
\end{figure*}

\subsection{Domestic events}
\label{sec:domestic}
Of the 922 EVENT articles, 604 reported a domestic event, 271 a foreign event and 47 could not be placed. The blocks differ in kind, not only in degree (Table~\ref{tab:domestic}). In the industrial block 97\% of EVENT articles were domestic, so the event share is the domestic event share (41.9\%). In the conventional block 44\% were domestic: ``mass shooting'' (1 of 53 EVENT articles domestic), ``drone attack'' (1 of 22) and ``explosive terror'' (15 of 48) are dominated by foreign coverage, while ``bomb threat'', ``terror-threat post'' and ``airport bomb report'' are almost entirely domestic. In the security block 18 of 95 EVENT articles were domestic (white-powder reports, suspicious mail, reactor trips and DPRK balloons; listed in the supplement), so the domestic event share is 2.0\%: one CBRN security article in 50 reports something that happened in Korea. On the human reference standard the figure is 2.5\% (Section~\ref{sec:human}), but only nine of these eighteen articles survive human coding, the reactor trips and DPRK balloon reports being among those the human coders placed outside the queries. Five further security EVENT articles report events in the DPRK; counting these as part of the domestic threat environment raises the share to 2.6\%. The 43 conventional EVENT articles with no codable place bound that block's domestic share between 18.1\% and 22.4\%. The domestic share is also a candidate explanation of the block contrast: CBRN queries retrieve foreign and geopolitical coverage that the other blocks do not, and the present design cannot separate that from the rarity of the hazards themselves.

\begin{table}[!htb]
\centering\small
\caption{Domestic and foreign events among EVENT articles in the gold set, by block. $p_E^{\mathrm{dom}}$ is the share of all sampled articles that report a domestic event.}
\label{tab:domestic}
\begin{tabular*}{\textwidth}{@{\extracolsep{\fill}}lrrrrrr@{}}
\toprule
Block & $n$ & EVENT & Domestic & Foreign & Unknown & $p_E^{\mathrm{dom}}$ (\%) \\
\midrule
Industrial & 967 & 418 & 405 & 10 & 3 & 41.9 \\
Security (CBRN) & 900 & 95 & 18 & 76 & 1 & 2.0 \\
Conventional terror & 1,000 & 409 & 181 & 185 & 43 & 18.1 \\
\bottomrule
\end{tabular*}
\end{table}

\subsection{Amplification}
\label{sec:amp}
Grouping EVENT articles into unique events, with events merged across the queries of a block, gives $m = 1.75$ articles per event in the industrial block, 1.12 in the security block and 1.19 in the conventional block, and amplification factors of 4.05 (95\% CI 3.82--4.30), 10.59 (9.06--12.57) and 2.90 (2.73--3.10); the pooled value is 4.29 (4.11--4.49) (Fig.~\ref{fig:decomp}b). Clustering within queries instead, so that an accident retrieved by two industrial queries counts twice, gives $m = 1.63$, 1.10 and 1.15 and $A = 3.76$, 10.47 and 2.81; the merge affects the industrial block most because its queries overlap. Query-level values (Fig.~\ref{fig:decomp}d) use within-query clustering. These sample-based values are lower bounds, and the block ordering of $m$ should not be read as a property of the hazards: $m$ rises with the sampling fraction (Spearman $\rho = 0.31$ across queries, within-query clustering; mean 1.11 for the eight queries sampled below 5\% and 1.36 for the sixteen sampled above 20\%), and the industrial block, whose queries are small, is the block in which duplicate reports are most likely to be captured. The block contrast in $A$ is carried by $p_E$, which is unbiased under the design but is a model label; on the human reference standard the security value rises to 22.5, and to 33--40 when the duplication factor is also put on the human-confirmed events (Section~\ref{sec:human}). The clustering rule itself is not the source of the ordering: varying the window between one and seven days moves $m$ by less than 0.05 in every block, and relaxing the place and object criteria raises $m$ in every block without changing their order (S11). Query-level values range from 1.6 (``bomb threat'') to 100 (``bioterrorism'', one EVENT article in 100); the extremes are determined by one to three EVENT articles and, because sampling fractions differ, query values are not comparable with one another.

\subsection{Reference bias}
\label{sec:refbias}
Reference bias is concentrated in rare hazards (Fig.~\ref{fig:refbias}a, Table~\ref{tab:anchor}). With title, excerpt and keywords, 89.3\% of the 224 Novichok articles mention a past Novichok poisoning: 79.0\% mention Navalny and 25.4\% mention Skripal or Salisbury. Of 41 sarin articles, 56.1\% mention Aum Shinrikyo, Tokyo or Ghouta. Of 2,704 radiation-leak articles, 44.9\% mention Fukushima (37.7\%) or Chernobyl (11.7\%). Of 309 anthrax articles, 38.2\% mention the United States letters. Of 363 hydrogen fluoride articles, 33.1\% mention Gumi, fourteen years after the leak. The three industrial controls scored on the same Gumi anchor return 2.1\% (chlorine), 2.4\% (sulfuric acid) and 1.4\% (nitric acid). On title and excerpt alone every value is lower, since the keyword field draws on the full text (Novichok 51.3\%, sarin 24.4\%, radiation leak 26.9\%, anthrax 14.6\%, hydrogen fluoride 10.5\%; controls 0.0--0.8\%), but the ordering and the control contrast are the same on either basis. Adding bare year tokens changes no value by more than 0.6 points.

Scoring the controls against hydrogen fluoride's anchor invites an obvious objection: a hazard with no canonical disaster of its own has no reason to mention someone else's, so the contrast could be a definition rather than a measurement. We therefore repeated the comparison with a rule that names no anchor at all. An article counts as citing a past occurrence if its title or excerpt contains a year between 1990 and 2020 within forty characters of an incident word (accident, disaster, leak, explosion, attack, exposure, poisoning). On that rule hydrogen fluoride coverage cites a past occurrence in 10.7\% of articles (95\% CI 7.7--14.1) against 1.7\% (0.4--3.3) for chlorine, 2.0\% (0.4--3.9) for sulfuric acid and 3.7\% (1.4--6.5) for nitric acid: ratios of 6.5, 5.5 and 2.9, with Fisher's exact $p < 0.005$ in all three comparisons. Reading the matched passages shows why the control values are not zero and also why they are not citations: in the three control queries the year tokens are administrative, giving the year a waterworks began operating or the span of a chemical-accident tally, whereas in the hydrogen fluoride query they name the 2012 Gumi leak. The controls are not declining to cite Gumi. They have no canonical occurrence to cite, which is the condition the contrast is meant to capture, and it is visible under a rule that was not written around any particular event. This anchor-free rule is a coarser instrument than the lexicons, since it catches only citations that carry a year, and it should not be compared with the lexicon values directly.

Two objections to this comparison can be answered from the same labels. The first is that $\rho$ and the event share are not independent, since a hazard with few events has more of its volume available for every non-event category, anchor citations included. Restricting to non-event articles removes that. These rates are computed on title and excerpt, the basis available for every gold-set article, and are therefore not comparable with the keyword-based values above; among gold-set articles the coders did not label EVENT, the anchor citation rate is 52.0\% for Novichok, 30.6\% for radiation leak, 21.1\% for sarin, 18.3\% for hydrogen fluoride and 9.7\% for anthrax, while on the same basis the three industrial controls return 0 of 123 such articles, an upper bound of 2.4\% by the rule of three. The contrast is at least fourfold on the strictest reading and is not an artefact of event scarcity. The second objection is that the controls are scored against an anchor belonging to a different substance, so what they measure is cross-hazard spillover rather than a general tendency to look backward. The comparison asks whether a hazard's own canonical disaster is cited at a rate that a neighbouring hazard's is not; it does not establish a backward-looking baseline for coverage in general. The security block has no within-block control (Section~\ref{sec:limits}).

\begin{table}[!htb]
\centering\small
\caption{Anchor mentions by query. $\rho_{\mathrm{te}}$ uses title and excerpt; $\rho_{\mathrm{kw}}$ adds the archive's keyword field (year tokens excluded in both). Gold-set columns use the keyword basis: articles mentioning the anchor ($k$ of $n$), the share of those coded EVENT, and the share of all sampled articles that mention the anchor and report no event.}
\label{tab:anchor}
\setlength{\tabcolsep}{3pt}
\begin{tabular*}{\textwidth}{@{\extracolsep{\fill}}llrrrrrr@{}}
\toprule
Query & Anchor lexicon & $V$ & $\rho_{\mathrm{te}}$ & $\rho_{\mathrm{kw}}$ & Gold $k/n$ & EVENT & Anchor, \\
 & & & (\%) & (\%) & & (\%) & no event (\%) \\
\midrule
Novichok & Navalny; Skripal/Salisbury & 224 & 51.3 & 89.3 & 90/100 & 28 & 65 \\
Sarin attack & Aum, Tokyo; Ghouta & 41 & 24.4 & 56.1 & 23/41 & 13 & 49 \\
Radiation leak & Fukushima; Chernobyl & 2,704 & 26.9 & 44.9 & 44/100 & 20 & 35 \\
Anthrax terror & US letters & 309 & 14.6 & 38.2 & 41/100 & 17 & 34 \\
HF leak & Gumi & 363 & 10.5 & 33.1 & 32/100 & 19 & 26 \\
Cl$_2$ gas leak (control) & Gumi & 241 & 0.8 & 2.1 & 0/100 & -- & 0 \\
H$_2$SO$_4$ spill (control) & Gumi & 254 & 0.0 & 2.4 & 3/100 & 0 & 3 \\
HNO$_3$ leak (control) & Gumi & 216 & 0.0 & 1.4 & 2/100 & 50 & 1 \\
\bottomrule
\end{tabular*}
\end{table}

An anchor lexicon can also match new developments of the anchor's own story, and the two largest values show this. Navalny mentions in Novichok coverage rise from 33\% of articles in 2021 to 97\% in 2024 (98 articles), the year of his death in custody. The coders treated the death, reported as a suspected poisoning, as an event of the query's object: 17 of the 25 Novichok EVENT articles in the gold set report it, and the rest report other suspected poisonings of Russian opponents; under a strict reading in which only a new use of Novichok counts, the Novichok event share is 0--1\% and its coverage is reference from end to end. The poisoning the 2024 coverage recalled was four years old. Fukushima mentions in radiation-leak coverage peak at 66\% in 2023 (894 articles), the year treated water was first discharged, which was itself an event. The gold-set cross-tabulation separates the two uses: among sampled articles that mention the anchor, the share coded EVENT is 28\% for Novichok (90 articles), 20\% for radiation leak (44), 19\% for hydrogen fluoride (32), 17\% for anthrax (41) and 13\% for sarin (23), so roughly four anchor mentions in five occur in articles that report no event. Expressed as a share of all articles, 65\% of Novichok coverage, 49\% of sarin coverage and 26\% of hydrogen fluoride coverage cites the past event without reporting a new one. The REF label, which requires the past event to be the article's main subject, captures only a small part of this (2.2--3.8\% of articles by block); reference bias operates mainly through insertion, as in the headline ``SK Hynix fluorine accident recalls the Gumi nightmare'', which reports a new event and cites the anchor in the same breath.

\begin{figure*}[!t]
\centering
\includegraphics[width=\textwidth]{Fig3_reference_bias_ground_truth.pdf}
\caption{Reference bias and the official registry comparison. (a) Share of articles citing a past anchor event for five CBRN or rare-hazard queries and three same-block industrial controls scored against the same anchor (Gumi 2012); for Novichok and radiation leak the dominant anchor is labelled and the remainder is the other anchor of the same hazard. (b) Article count versus registry incident count for six substances (9 Sep 2021 to 28 Apr 2025), with the log--log regression of articles on incidents. (c) Articles per official incident with article and incident counts (descriptive only; the ratio is not used for inference).}
\label{fig:refbias}
\end{figure*}

\subsection{Comparison with the official registry}
Within the 3.6-year overlap window the registry records 366 chemical accidents, of which 142 name at least one of the six substances; counted by substance rather than by incident, and so counting nine mixed-acid incidents twice, the figure is 151. Articles per official incident range from 3.0 (sulfuric acid, 164 articles for 55 incidents) and 3.1 (nitric acid, 110 for 36) to 16.2 (chlorine, 162 for 10) and 20.3 (hydrogen fluoride, 244 for 12) (Fig.~\ref{fig:refbias}b,c). The article counts themselves are unrelated to the incident counts: the log--log slope of articles on incidents is $-0.07$ (95\% CI $-0.89$ to $0.76$, $r = -0.11$, $p = 0.83$; Spearman $\rho = 0.03$, $n = 6$). A five-fold difference in how often these substances are released produces no detectable difference in how many articles the corresponding queries return, and the query for hydrogen fluoride, with twelve incidents behind it, returned more articles than the query for sulfuric acid with fifty-five. The claim is about queries rather than substances, and two features of the design keep it there. The six queries are not matched on length, and because the archive treats a space as a logical AND, the three-token queries (ammonia, chlorine, sulfuric acid) are stricter than the two-token ones; the two substances with the most registry incidents are among the stricter queries. A substance may also be named differently across contexts, as hydrogen fluoride is in Korean industrial and trade reporting, so a single-spelling query recovers a different fraction for different substances. Neither can be corrected without recollecting the corpus under matched queries, which we did not do. What the registry matching itself contributes to the null is smaller: the substance fields are free text, and the original rule also catches nitrate salts, a toluene diisocyanate release, a stand-alone hypochlorite spill and fluoroboric acid. Recoding strictly, so that only releases of the named substance count, moves the counts to 11, 24, 8, 55, 13 and 31 and the slope from $-0.07$ to $-0.04$ ($p = 0.90$); the null is not an artefact of the matching. With six substances the interval is wide and the result is a null, not an estimate of a small effect. The hypothesis that the practice assumes, a slope of 1, is rejected ($t(4) = 3.6$, $p = 0.023$), and leaving out any one substance moves the slope only between $-0.19$ and $0.06$. The registry is not ground truth either. It is a second reporting instrument, with a statutory notification threshold and operator self-report, so its own probability of capturing an incident is unknown and need not be constant across substances; the comparison is therefore between two instruments that diverge. Measurement error in the registry counts would attenuate the slope towards zero, which is one reason not to read the point estimate as evidence of a negative relation, but attenuation alone cannot generate a negative estimate from a true slope of 1. It can account for the interval: a true slope of 1 attenuated by a reliability of 0.76, meaning that a quarter of the variance in log incident counts is error, would land at the interval's upper limit. Adjusting for severity by adding log casualties leaves three residual degrees of freedom on six substances. In that model the incident slope is 0.64 and the casualty slope is $-0.67$, so the fit buys a higher incident coefficient by assigning fewer articles to substances that injure more people. A severity-adjusted model is not identified at this sample size, and we report it only to record that the null in the direct model is not produced by ignoring severity. The two broad queries ``hazardous chemical release'' and ``toxic-gas evacuation'' return fewer articles than there are registry incidents (0.17 and 0.29 per incident) because they capture only a fraction of the 300 leak-type incidents; this is a recall limit of the query, not evidence about amplification.

The framework's own event estimates for the six substances, $\hat N = V p_E^{\mathrm{corpus}} d / m$, range from 25 (toluene) to 75 (hydrogen fluoride) domestic events against 14 and 12 registry incidents, and from 37 (nitric acid) to 61 (sulfuric acid) against 36 and 55. These use the sample $m$ and are therefore upper bounds; they exceed the registry for every substance, by a factor of 1.0 to 6.3. With the corpus $m$ of the industrial block (2.25) the estimates fall to 17--57 and two substances (sulfuric and nitric acid) fall below their registry counts. The ranking of the six substances by news-derived events is unrelated to their ranking by registry incidents (Spearman $\rho = -0.09$), so the decomposition removes the non-event coverage but does not by itself recover the registry's ordering. The most likely reason is a difference in what counts as an incident, since the registry records accidents that meet a statutory reporting threshold whereas the news also reports minor releases and false alarms, but we did not link individual news events to registry records and cannot separate that explanation from coding error.

\subsection{Classifier}
\label{sec:clf}
Table~\ref{tab:models} gives the full comparison; three results in it matter here. The keyword rule that underlies most volume-based practice reaches EVENT-$F_1$ 0.504 with precision 0.50, so half of what a keyword count calls an event is not one, and the proposed model beats it by 0.355 (Nadeau--Bengio corrected $t = 25.2$, $p < 0.001$). Adding the archive's incident taxonomy does not improve the model, by 0.014 macro-$F_1$ in the wrong direction, which the corrected test does not distinguish from zero ($p = 0.080$). And the model is weakest exactly where the decomposition matters most: per-block EVENT-$F_1$ is 0.911 (industrial), 0.867 (conventional) and 0.507 (security), the last resting on 95 EVENT articles among 900. The pre-specified sanity checks are met with one qualification: macro-$F_1$ rises with training fraction on the seed mean and on eight of ten seeds individually, the other two falling by 0.007 and 0.019 between the 75\% and 100\% fractions (Fig.~\ref{fig:clf}b); random-baseline accuracy is 0.200 as required; and block divergence stays within the specified factor of 1.5 on macro-$F_1$ (ratio 1.25) but not on EVENT-$F_1$ (1.8), for which it was not specified. Per-class sensitivity and false-positive rates against coder A's labels are in the supplement (S12); they are not the rates used in the misclassification correction, which come from the human standard of Section~\ref{sec:human}.

Leave-one-query-out evaluation, which withholds each query's vocabulary from training, gives pooled EVENT-$F_1$ of 0.85 for industrial queries, 0.79 for conventional queries and 0.38 for security queries (0.68 overall; per-query values in the supplement range from 0.00 to 1.00, with ``drone attack'' at 0.57 and ``DPRK nuclear test'' at 0.00), and predicted event shares of 46\%, 42\% and 6\% against true shares of 43\%, 41\% and 11\%. For the two control blocks the model transfers to unseen hazards with a loss of a few points; for CBRN queries it does not, which is why the corpus estimates below rest on post-stratification rather than classification for that block.

\begin{table}[!htb]
\centering\small
\caption{Classifier performance, 10 seeds $\times$ stratified 5-fold CV (mean $\pm$ SD over 50 folds). Proposed: character $n$-gram (2--4) TF--IDF with class-weighted logistic regression. The majority-class baseline predicts EVENT for every article; the keyword rule predicts EVENT when all query tokens are present and OTHER otherwise.}
\label{tab:models}
\setlength{\tabcolsep}{5pt}
\begin{tabular*}{\textwidth}{@{\extracolsep{\fill}}lccc@{}}
\toprule
Model & Accuracy & Macro-$F_1$ & EVENT-$F_1$ \\
\midrule
Majority class (predict EVENT) & 0.322 $\pm$ 0.001 & 0.097 $\pm$ 0.000 & 0.487 $\pm$ 0.001 \\
Random ($1/K$) & 0.200 $\pm$ 0.016 & 0.182 $\pm$ 0.015 & 0.246 $\pm$ 0.032 \\
Keyword rule & 0.322 $\pm$ 0.015 & 0.173 $\pm$ 0.008 & 0.504 $\pm$ 0.025 \\
Char $n$-gram, title only & 0.723 $\pm$ 0.019 & 0.648 $\pm$ 0.029 & 0.830 $\pm$ 0.018 \\
Word $n$-gram (1--2), title + excerpt & 0.745 $\pm$ 0.015 & 0.672 $\pm$ 0.026 & 0.854 $\pm$ 0.015 \\
Char $n$-gram + archive incident tag & 0.758 $\pm$ 0.016 & 0.700 $\pm$ 0.024 & 0.854 $\pm$ 0.015 \\
\textbf{Char $n$-gram, title + excerpt (proposed)} & 0.769 $\pm$ 0.015 & 0.714 $\pm$ 0.021 & 0.859 $\pm$ 0.015 \\
\quad industrial block & -- & 0.754 $\pm$ 0.047 & 0.911 $\pm$ 0.018 \\
\quad security block & -- & 0.602 $\pm$ 0.060 & 0.507 $\pm$ 0.098 \\
\quad conventional block & -- & 0.693 $\pm$ 0.038 & 0.867 $\pm$ 0.024 \\
\midrule
\multicolumn{4}{@{}l}{\footnotesize Proposed vs.\ keyword rule, EVENT-$F_1$: $\Delta = +0.355$ [0.352, 0.358], Nadeau--Bengio $t = 25.2$, $p < 0.001$.} \\
\multicolumn{4}{@{}l}{\footnotesize Archive incident tag: $\Delta$ macro-$F_1$ $= -0.014$ [$-0.018$, $-0.010$], corrected $p = 0.080$; $\Delta$ EVENT-$F_1$ $= -0.005$, $p = 0.35$.} \\
\bottomrule
\end{tabular*}
\end{table}

Inter-coder agreement on the 600 double-coded articles is $\kappa = 0.790$ (95\% CI 0.751--0.829); for the binary EVENT-versus-rest decision it is 0.875 (Fig.~\ref{fig:clf}c), but by block it is 0.90 (industrial), 0.92 (conventional) and 0.55 (security), and coder B's event share on the subset is 8.8\% in the security block against coder A's 9.9\% (31.8\% against 33.3\% overall). The security block is where the two coders, like the classifier, agree least. Of 94 disagreements, 27 were between OTHER and IRRELEVANT, 24 between INST and OTHER, 16 between EVENT and IRRELEVANT and 13 between EVENT and OTHER; the EVENT boundary that determines $p_E$ is the most stable one.

\begin{figure*}[!t]
\centering
\includegraphics[width=\textwidth]{Fig4_classifier_validation.pdf}
\caption{Classifier performance and label validation. (a) Macro-$F_1$ and EVENT-$F_1$ for baselines, ablations and the proposed model (10 seeds by stratified 5-fold CV, mean $\pm$ SD). (b) Learning curve over training fraction (one held-out fold per seed). (c) Confusion matrix between the two model coders, A and B, on the 600-article double-coded subset (row-normalised; $\kappa = 0.790$, EVENT-versus-rest $\kappa = 0.875$). (d) Held-out performance by block; the security block contains 95 EVENT articles.}
\label{fig:clf}
\end{figure*}

The archive's incident taxonomy agrees weakly with the EVENT label ($\phi = 0.108$, $\kappa = 0.104$, odds ratio 1.60, 95\% CI 1.36--1.87). Among articles tagged ``industrial accident: explosion'' 42\% reported an event and 34\% reported drills or inspections; among those tagged ``nuclear accident'' 10\% reported an event. The taxonomy encodes topic, not occurrence, which is why it does not help the classifier.

\subsection{Human validation}
\label{sec:human}
Two authors coded the 400-article subset independently under the same codebook, blind to the model labels and to each other. Agreement between them is $\kappa = 0.848$ (95\% CI 0.804--0.892) over the five components and 0.905 (0.853--0.958) for the EVENT-versus-rest decision, with 90.0\% raw agreement; by block the five-class value is 0.930 (industrial), 0.817 (conventional) and 0.810 (security). Forty articles are coded differently. Twenty-five of these lie between IRRELEVANT and either INST or OTHER, but twelve touch the EVENT boundary, and nine of the twelve are the same query: articles retrieved by ``radiation leak'' that one coder read as an event and the other as background. That is a rule disagreement rather than noise, and it is the largest single source of uncertainty in the security estimate below. The 360 articles on which the two coders agree are used as the reference standard, and the effect of that choice is quantified.

Against the standard, coder A reaches $\kappa = 0.396$ (0.338--0.454) over the five components and 0.499 (0.410--0.587) on the EVENT boundary. The block pattern matters more than the pooled value: 0.938 in the conventional block and 0.689 in the industrial block, against 0.236 (0.125--0.348) in the security block. Coder B, on the 66 subset articles it also coded, reaches 0.810 on the EVENT boundary.

Table~\ref{tab:kappa} sets these against the model--model and human--human values on matching metrics, since a five-class coefficient and a binary one are not comparable and neither is a pooled value against a block value. Read down the EVENT-boundary column, the two models agree with each other at 0.875 and with the human standard at 0.499; within the security block the same two figures are 0.546 and 0.236, while the two human coders agree with each other there at 0.793 (0.670--0.916). The human standard is therefore reliable but not perfect in the block that carries the claim, and the model--human gap is wider than the shortfall in the standard's own reliability can account for.

\begin{table}[!htb]
\centering\small
\caption{Agreement on matching metrics. Each cell is Cohen's $\kappa$ with an asymptotic 95\% interval, computed on the articles the two coders in that row share: 600 for the two models, 400 for the two humans, 360 for the human standard against coder A and 66 against coder B. Five-class and binary coefficients are not comparable with each other, and neither is a pooled value with a block value.}
\label{tab:kappa}
\footnotesize\setlength{\tabcolsep}{4pt}
\begin{tabular*}{\textwidth}{@{\extracolsep{\fill}}lrrrrr@{}}
\toprule
Pair & $n$ & \multicolumn{2}{c}{All blocks} & \multicolumn{2}{c}{Security block} \\
\cmidrule(lr){3-4}\cmidrule(lr){5-6}
 & & five-class & EVENT boundary & five-class & EVENT boundary \\
\midrule
Model A vs model B & 600 & 0.790 & 0.875 & 0.671 & 0.546 \\
Human H1 vs human H2 & 400 & 0.848 & 0.905 & 0.810 & 0.793 \\
Human standard vs model A & 360 & 0.396 & 0.499 & 0.284 & 0.236 \\
Human standard vs model B & 66 & 0.551 & 0.810 & 0.491 & 0.553 \\
\bottomrule
\end{tabular*}
\end{table}

The security-block disagreement is systematic. Of the 79 security articles in the reference standard that coder A called EVENT, 61 were not events to either human coder, and 56 of those were coded IRRELEVANT: the article reports something that happened, but not of the kind the query names. The largest groups are the Novichok query (24), radiation leak (13), DPRK nuclear test (5) and anthrax (5) (Fig.~\ref{fig:human}b); characteristic cases are the death of Alexei Navalny in prison retrieved by the Novichok query, a ballistic-missile launch retrieved by the nuclear-test query, and a train derailment releasing toxic gas retrieved by the chemical-weapons query. The codebook's means-mismatch rule (Section~\ref{sec:coding}) puts these outside the query; the model, working from the same text, treated topical adjacency as sufficient. The human coders did not add this rule; it is in the version of the codebook both model coders were given, with the nuclear-for-chemical case written out as a worked example. That covers the clearest group, such as a ballistic-missile launch retrieved by the nuclear-test query. It does not cover all of them. Whether a death by suspected poisoning is an occurrence of the Novichok query's object, or a reactor trip an occurrence of the radiation-leak query's, is a boundary the codebook states in principle but does not settle with an example, and those are the two queries on which the human coders themselves disagreed most (nine of the twelve EVENT-boundary disagreements are the radiation-leak query). The divergence therefore has two sources: a shared instruction the models did not apply, and a boundary the instruction left underspecified, which the models resolved permissively and the human coders strictly. Neither source is visible from model--model agreement. The Novichok query is where this cuts finest: of the 25 articles coder A called events there, the human coders kept one, a report that the poisoning agent was Novichok, and rejected the other 24, which report the death without the agent. The distinction they drew is between an article about a use of the query's object and an article about its consequences. In the opposite direction only 4 human EVENT articles were missed. The full table is 18 true positives, 61 false positives, 4 false negatives and 136 true negatives, so coder A's precision on the EVENT class in this block is 0.23: fewer than one in four of its event calls survives human coding, which is the number a practitioner reusing a model coder needs and is more informative here than $\kappa$, whose value is depressed by the low prevalence. The 4 false negatives sit in the stratum sampled at 0.19 and therefore enter the weighted estimate at 5.4 each, which is why the sensitivity interval is as wide as 0.28--0.83; adding or removing one of them moves the security event share from 5.0\% to 5.6\% or 4.3\%, inside the range already reported for coder disagreement. Coder A therefore over-detects events in the security block, and the over-detection is concentrated in the rare hazards that carry the highest reference bias, where a past or adjacent occurrence is what the coverage is about.

Table~\ref{tab:human} and Fig.~\ref{fig:human} give the consequences. Post-stratified onto the known cell sizes, the human event share is 38.0\% (95\% CI 33.6--40.9) in the conventional block, 33.8\% (25.7--41.9) in the industrial block and 5.0\% (2.5--7.9) in the security block, against coder A's 40.9\%, 43.2\% and 10.6\%. The security share is halved; the control blocks move by less than their intervals. Conditioning on topical relevance, as in Section~\ref{sec:composition}, gives human event shares of 52.1\% (42.0--63.0) industrial, 59.6\% (51.3--69.5) conventional and 13.9\% (7.3--21.8) security. The security interval does not reach either control's, so on human labels as on model labels the block contrast is not produced by the queries retrieving with different precision; these intervals come from a within-cell bootstrap, and the query-clustered version is given at the end of this section. Coder A's sensitivity and false-positive rate against the standard are, in the same order, 1.00 and 0.05 (conventional), 0.92 and 0.18 (industrial), and 0.48 and 0.09 (security). Applying the Rogan--Gladen correction to coder A's gold-set share with these rates returns the post-stratified estimate exactly (largest absolute difference across the three blocks $5.6 \times 10^{-17}$), because under a common weighting the two are the same estimator written differently, so the agreement between them is not corroboration. The correction is informative only when it is carried to a quantity estimated outside the subset, which is done for the corpus in Section~\ref{sec:corpus}.

The security estimate depends on how the 40 human disagreements are treated, and the dependence is bounded. Five treatments produce four distinct values: 5.0\% on consensus, 4.4\% taking H1 alone or resolving every disagreement as a non-event, and 6.0\% taking H2 alone or counting either coder's EVENT. The defensible range is therefore 4.4--6.0\%, and it is wider than the bootstrap interval suggests on one side only, since the disagreements move the estimate in a single direction. A sixth treatment, in which all 40 disagreements are coded as events, gives 12.1\%, but it is not a bound on anything the coders said: 28 of the 40 disagreements lie between non-event categories, so that treatment promotes to EVENT articles neither coder called an event.\footnote{It is the only treatment that reaches coder A's 10.6\%, which is worth recording because it shows how far the model's figure sits from any reading of the human coding: it is recovered only by a rule that overrides both coders.} Most of the remaining uncertainty is one query. Dropping ``radiation leak'', which supplies nine of the twelve EVENT-boundary disagreements, leaves the consensus estimate at 5.0\% and the defensible range at 4.8--5.4\%, so the point estimate is insensitive to the query and the width is largely its rule ambiguity.

Recomputing amplification on the human event share, holding the duplication factor at its model-based value from Table~\ref{tab:components}, gives $A = 22.5$ (14.2--44.0) for the security block, 5.2 (4.2--6.8) for the industrial block and 3.1 (2.9--3.5) for the conventional block. These are hybrid quantities, with a human numerator share and a model-based duplication factor, and the duplication factor is the weaker half: $m = 1.12$ was computed over the 95 articles coder A called events in this block, of which 18 survive the human standard. Because the security block is small enough to settle by hand, one of us assigned event date, place and object to all 22 articles the two human coders agreed were events and passed them through the same clustering function, a third coding pass by a single analyst who was not blind to the analysis. On those events the duplication factor is between 1.45 and 2.00 depending on the treatment of one large cluster and on whether the articles are weighted by their selection probability, against 1.12 on coder A's labels, so $A$ would be 29 to 40 rather than 22.5. We keep 22.5 as the reported figure: it is the conservative end, it does not rest on a single non-blind pass, and correcting the denominator moves the estimate away from the controls rather than towards them. The coding, the clusters and the variants are in the supplement (S10) and the deposit. The block ordering is unchanged and the gap between the security block and the controls is wider than the model labels showed.

The domestic event share is also revised, though it moves little. All 18 security articles that coder A labelled EVENT and placed in Korea fall in the cell that was sampled in full, so all were human-coded; 9 survive as events to both coders, the rejected 9 being reactor trips and DPRK balloon reports that the human coders placed outside the queries. Against this the human coders identified domestic events coder A had missed, drawn from the cell sampled at 0.19. Post-stratified on the same consensus basis as the event shares above, the domestic share of security coverage is 2.5\% (0.8--4.6) on human labels against 2.0\% on coder A's; counting the 40 disagreements as non-events gives 2.2\%. The figure is stable but its composition is not: it holds at about one article in forty on the human standard against one in fifty on coder A's through two errors of opposite sign, so the agreement is not a validated count.

The validation covers only part of the results. It licenses the event shares of this section, the amplification values recomputed from them, and the corpus correction of Section~\ref{sec:corpus}. It does not license three things reported on coder A's labels elsewhere: the duplication factors, the gold-set cross-tabulation of Table~\ref{tab:anchor}, and the classifier of Section~\ref{sec:clf}. The query-level block tests of Section~\ref{sec:composition} are re-run on the human labels below. The classifier is the sharpest case. It is trained and scored on those labels, so its security-block EVENT-$F_1$ of 0.51 measures agreement with an annotator now known to be wrong in that block, not performance on the construct. The block contrast survives correction. On the ratio scale it widens: industrial to security rises from 4.1 to 6.8 and conventional to security from 3.9 to 7.7. On the percentage-point scale the security-to-industrial gap narrows from 32.6 to 28.8, because the industrial share also falls.

The contrast can also be tested rather than described. The block comparisons of Section~\ref{sec:composition} were run on model labels with the query as the unit; we repeat both on the human labels. Two things must be separated here. The first is the interval on each block's share. Section~\ref{sec:stats} treats the query as the unit for between-block comparisons, but the human intervals given above come from a within-cell bootstrap, which ignores that clustering. Resampling whole queries within each block instead, and post-stratifying each resample onto the same block cells, gives 5.0\% (0.6--10.9) for the security block against 33.8\% (22.9--39.9) industrial and 38.0\% (30.2--40.9) conventional; conditioning on topical relevance gives 13.9\% (1.9--31.2) against 52.1\% (36.4--72.0) and 59.6\% (52.4--68.4). The security interval lies clear of both controls on both bases, so the contrast is not produced by within-query clustering, and this is the form in which we state it.

The second is the rank test itself, and here the subset is thin. The human sample was designed to estimate block shares, not to power a per-query test in the control blocks: it holds 219 articles across the ten security queries but a median of 7 across each control block's ten, with a minimum of 4, and three queries have no agreed article in one of their two cells. On the human labels the omnibus test still rejects for the event share ($H(2) = 6.37$, $p = 0.041$, $\eta^2 = 0.16$, $\epsilon^2 = 0.22$; medians 36\%, 1\% and 44\%), although $\eta^2$ here sits only just above the 95th percentile of its permutation null (0.15) and its resampling interval includes zero, and the security block lies below the conventional block (Mann--Whitney $U = 18$, Holm-adjusted $p = 0.047$, $r = 0.64$); the security-to-industrial comparison is in the same direction but does not reach significance ($U = 26$, $p = 0.139$, $r = 0.48$), and the two controls remain indistinguishable ($p = 0.85$). Conditioning on topical relevance, the rank test does not reject at all ($H(2) = 3.24$, $p = 0.198$): with five to eleven agreed articles per control query the conditional share is 0 or 1 for several of them and the test is reduced to a handful of tied ranks. We therefore do not claim a query-level rank test of the conditional contrast on human labels. What the human labels do support is the block-level statement with query-clustered intervals, the omnibus test on the event share, and the security-to-conventional pair. The same tests on model labels reject more strongly in every case ($H = 10.06$ and $9.29$), but that is a difference in how many articles stand behind each query rather than a difference in the labels: the model coded 100 articles in every query and the human subset has a median of 7 in the control blocks. The direction and the ordering are the same on both label sets. The reference-bias measurements, the registry comparison and the corpus volume series use no component labels and are untouched by any of this.

\begin{figure*}[!t]
\centering
\includegraphics[width=\textwidth]{Fig5_human_validation.pdf}
\caption{Human validation of the model labels. (a) Event share by block: coder A's own labels against the post-stratified human reference standard with its 95\% bootstrap interval, and the range spanned by the five treatments of the 40 human disagreements. $\kappa$ is human-versus-coder-A agreement on the EVENT boundary for that block. Only the security block moves beyond its interval. (b) The 61 security articles coder A called EVENT that both human coders rejected, by query. Fifty-six were coded IRRELEVANT under the codebook's means-mismatch rule: the article reports an occurrence, but not of the kind the query names.}
\label{fig:human}
\end{figure*}

\begin{table*}[!t]
\centering\small
\caption{Human validation of the model labels on the 400-article subset. $n$ is articles drawn and, in parentheses, retained after conditioning on agreement between the two human coders. TPR and FPR are coder A's sensitivity and false-positive rate against that standard under the same post-stratification. Human $p_E$ is the post-stratified event share; the range column spans the treatments of the 40 disagreements that follow one or both coders (each coder alone, either coder's EVENT, and all disagreements as non-events); the further treatment that codes every disagreement as an event is excluded, since 28 of the 40 lie between non-event categories and neither coder called them events. $A$ is amplification recomputed on the human event share with the duplication factor held at its model-based value, and is a lower bound. The conventional upper limit coincides with coder A's share because no human event was found in that block's non-EVENT cell, which caps the estimator at coder A's value; the coincidence is structural and not evidence that the two agree.}
\label{tab:human}
\footnotesize\setlength{\tabcolsep}{3pt}
\begin{tabular*}{\textwidth}{@{\extracolsep{\fill}}lrrrrrrr@{}}
\toprule
Block & $n$ & TPR & FPR & \multicolumn{3}{c}{Gold-set $p_E$ (\%)} & $A$ (human) \\
\cmidrule(lr){5-7}
 & & & & coder A & human & range & \\
\midrule
Industrial & 78 (74) & 0.92 & 0.18 & 43.2 & 33.8 [25.7, 41.9] & 33.5--33.8 & 5.2 [4.2, 6.8] \\
Security & 245 (219) & 0.48 & 0.09 & 10.6 & 5.0 [2.5, 7.9] & 4.4--6.0 & 22.5 [14.2, 44.0] \\
Conventional & 77 (67) & 1.00 & 0.05 & 40.9 & 38.0 [33.6, 40.9] & 35.4--38.2 & 3.1 [2.9, 3.5] \\
\bottomrule
\end{tabular*}
\end{table*}

\subsection{Full-corpus estimates}
\label{sec:corpus}
Post-stratifying the gold labels to the corpus gives an event share of 27.1\% (within-stratum 95\% CI 24.9--29.5) in the industrial block, 9.0\% (5.2--13.2) in the security block, 31.8\% (26.2--37.5) in the conventional block and 25.4\% (21.7--29.4) overall (Table~\ref{tab:fullcorpus}). Resampling whole queries widens these intervals considerably (industrial 11--55\%, security 5--23\%, conventional 24--56\%), because with ten queries per block the corpus share depends on which large queries happen to be in the block; the security estimate in particular is set by ``DPRK nuclear test'' (23,291 articles, 15,227 of them in 2022, with 7 EVENT articles in its 100 sampled). The corpus values are lower than the gold-set values because the gold set weights every query equally whereas the corpus is dominated by its largest queries: ``DPRK nuclear test'' is 78\% of the security block (23,291 articles, 15,227 of them in 2022) and ``drone attack'' is 65\% of the conventional block (48,812 articles). Excluding those two queries raises the block shares to 15.4\% and 50.0\%. The text-rule duplication factor on the corpus is 2.25 and 2.59 for the industrial and conventional blocks, larger than the sample values as expected when duplicates are no longer lost to sampling. We do not convert these into corpus-level amplification factors: the rule that produces them agrees with the coders on 36\% of article pairs, and a ratio built on it would carry that error into the quantity a reader is most likely to reuse. For the security block the classifier predicts EVENT for 2.7\% of articles against a post-stratified share of 9.0\%, and at its false-positive rate of 2.0\% most of those predictions are expected to be wrong, so a duplication factor computed on them is not an estimate of anything; the security corpus $A$ is therefore not reported. Correcting the post-stratified shares for coder A's measured misclassification (Section~\ref{sec:human}) gives 28.4\% (20.6--35.4) in the conventional block and 11.9\% (0.5--21.4) in the industrial block, with uncertainty in the sensitivity and false-positive rate propagated alongside uncertainty in the observed share. In the security block the corrected interval runs to an upper limit of 12.1\% with a lower limit at the boundary of the admissible range; we report that bound and not the point estimate it brackets, because the observed corpus share sits close to the false-positive rate and the point estimate is an artefact of that proximity. Unlike the gold-set correction, this one is not a restatement of the human estimate: it applies rates measured on the subset to a share estimated from the whole corpus. Only the conventional value is informative. In the industrial block the false-positive rate is itself uncertain (0.09--0.26) and the corrected interval spans almost the whole admissible range; in the security block the observed corpus share of 9.0\% sits close to the false-positive rate of 0.09, which is why only the bound is reported. The correction is only as precise as the rates it uses, and those rates come from 400 articles. The corrected interval therefore gives an upper bound on the corpus security event share, and the data supply no lower bound. The gold-set human estimate of 5.0\% is the measurement, and the corpus figure of 9.0\% should be treated as the model-label version of the same quantity, not as a separate finding.

\begin{table}[!htb]
\centering\small
\caption{Full-corpus estimates ($V = 109{,}610$). $p_E$ by post-stratification of gold labels over 178 query-by-year strata, with the within-stratum 95\% CI and, in brackets, the query-resampling CI; $m$ from the text clustering rule applied to classifier-predicted EVENT articles, indicative only: the rule is tuned on the gold-set clusters it is then compared with, agrees with the coders on 36\% of article pairs, and its gold-set $m$ moves between 1.24 and 1.59 across the six settings (Section~\ref{sec:stats}). n/e, not estimable: the classifier's predicted event share (2.7\%) is within sampling error of its false-positive rate. No corpus-level $A$ is reported: it would divide a design-based $p_E$ by a duplication factor whose pairwise agreement is 0.36, and the sample-based $A$ of Section~\ref{sec:amp} is the estimate we stand behind.}
\label{tab:fullcorpus}
\setlength{\tabcolsep}{3pt}
\begin{tabular*}{\textwidth}{@{\extracolsep{\fill}}lrrllr@{}}
\toprule
Block & $V$ & Pred.\ EVENT & $p_E^{\mathrm{corpus}}$ \% [CI] & query CI & $m^{\mathrm{corpus}}$ \\
\midrule
Industrial & 4,633 & 1,254 & 27.1 [24.9, 29.5] & [11.4, 55.3] & 2.25 \\
Security (CBRN) & 29,688 & 816 & 9.0 [5.2, 13.2] & [4.6, 22.8] & n/e \\
Conventional terror & 75,289 & 21,751 & 31.8 [26.2, 37.5] & [23.7, 55.7] & 2.59 \\
All & 109,610 & 23,821 & 25.4 [21.7, 29.4] & [16.5, 44.8] & -- \\
\bottomrule
\end{tabular*}
\end{table}

The ``drone attack'' query also illustrates why a volume trend is not an incidence trend. Its monthly rate rose from 311 articles in 2022 to 2,322 in the first eight months of 2026, a 7.5-fold increase (Fig.~\ref{fig:framework}c); 21 of its 22 sampled events were foreign, almost all war reporting. No change in domestic incidence is involved, and a planner reading the trend as a domestic signal would be reading the wrong country.

\subsection{Differential capture as a competing explanation}
\label{sec:capture}
The foreign share supplies a rival explanation for the low CBRN event share. A foreign occurrence must clear a higher threshold of consequence, proximity or political salience before a Korean outlet reports it at all \cite{galtung1965,vanbelle2000,yan2015,zhu2018}, and 76 of the 95 security EVENT articles are foreign. A low security share could therefore mean that CBRN events are rare, or that the archive cannot see the ordinary ones abroad; the same figure would be obtained in a country where CBRN incidents were common but unreported, and the two readings are not separable in these data.

Two features of the results bear on which dominates. The conventional block gives a within-hazard comparison: its queries retrieve both domestic and foreign occurrences, and its event share is high in both, which suggests that foreignness alone does not suppress event coverage once a hazard is newsworthy at all. The industrial block, whose events are almost entirely domestic, has an event share indistinguishable from the conventional block's, so the domestic--foreign axis does not order the blocks. Foreignness is therefore not what separates the security block, whose coverage is dominated by commentary, preparedness and retrieval error. Differential capture remains a real limit on the security estimate, and it is why the domestic event share is reported separately rather than folded into the block figure. The event-data literature has developed capture-recapture estimators for exactly this problem \cite{hendrix2015,weidmann2016,weidmann2014,dietrich2020}; applying them here would need a second independent source of CBRN incident reports, which for this hazard class is what does not exist.

\subsection{Sensitivity to query assignment}
\label{sec:sens}
The block contrast does not depend on the queries a reader might assign differently, nor on the way the queries are worded. Excluding ``DPRK nuclear test'', ``chemical accident response'', ``reactor radiation exposure'' or ``drone attack'', singly or all four together, leaves the Kruskal--Wallis $p$ at or below 0.01 and the security block below both controls at Holm-adjusted $p \le 0.03$; moving ``reactor radiation exposure'' into the security block instead of excluding it strengthens every contrast ($p = 0.0013$). Half the security queries name a threat without an occurrence noun, which is a candidate confound; the five that do contain one still fall below the industrial block ($p = 0.043$), so wording explains part of the within-block spread but not the contrast. The full sensitivity analysis, the per-query occurrence-noun coding and the by-year shares are given in the supplement (S9).

\section{Discussion}
\label{sec:discussion}

\subsection{What the decomposition changes for practice}
The immediate use of Eq.~\eqref{eq:A} is arithmetic. A planner who reads 1,000 articles about a CBRN hazard as 1,000 signals is, on the human-anchored estimate, reading at most 23 to 71 events, and on the model labels at most 80 to 110; for an industrial substance the same volume corresponds to at most 147--239 events, almost all domestic, and for a conventional threat to at most 283--345, fewer than half of them domestic. The two coding routes differ by a factor of two for CBRN hazards, by nine percentage points for industrial substances and within sampling error for conventional threats. That ordering is itself part of the finding: the rarer the hazard, the less reliable automated coding of its coverage was. It was observed for two models from one vendor, one prompt, one language and one codebook version, with rarity confounded against foreignness and query wording, so it is a hypothesis for replication rather than a property of model annotation. The three-block design does show that the low event share belongs to hazards that are rare, feared and institutionally active rather than to news in general. Industrial accidents and conventional terrorism, which differ in almost every other respect, have event shares this design cannot tell apart: the difference is 1.2 percentage points with a 95\% interval from $-17.5$ to $19.9$, and at ten queries per block the smallest difference detectable with 80\% power is 28 points. That is a statement about power, not evidence of equality. CBRN security hazards stand apart from both by a margin the same design resolves easily.

The composition tells the planner where the excess comes from. For CBRN hazards it comes from the near-absence of events beneath a large volume of preparedness coverage, commentary and off-topic retrieval, rather than from repetition of real events. Drills, surveillance notices and legislative activity are the institutional response the planner is part of, not noise, but they are the system talking about itself, and counting them as threat signal creates a loop in which preparedness activity generates coverage that justifies more preparedness activity. This loop is closer to the self-reinforcing news wave \cite{vasterman2005,wien2009} than to the amplification stations of the social amplification framework \cite{kasperson1988,pidgeon2003,kasperson2022}, whose quantitative applications relate coverage to societal impact \cite{renn1992,burns1993}; $p_E$ and the INST share make the loop a number for a given hazard. The same loop is visible in anthrax coverage: responders are called to suspicious powders that usually turn out to be benign, which is why field triage protocols for white powders exist \cite{kim2026whitepowder}, and each call generates coverage that a keyword count would score as an anthrax event.

The registry comparison is a null result with a practical edge. Across six substances whose incident counts differ five-fold, article counts do not differ, so hydrogen fluoride with twelve incidents drew more coverage than sulfuric acid with fifty-five. A budget allocated on coverage would, in this case, fund the substance that appears least often in the registry, and part of the reason is the reference load itself: a third of hydrogen fluoride coverage cites Gumi. The comparison cannot settle which instrument should govern the allocation; it shows that an agency using coverage is not using a proxy for the registry. The framework's own event estimates, as upper bounds, exceed the registry for every substance and do not recover its ordering, which sets a limit on what news can be asked to do: it can be cleaned of non-event coverage, but it cannot be made to apply the registry's definition of an incident, and a record-level linkage of news events to registry incidents, which we did not perform, is the way to find out how far the two definitions diverge. The finding is consistent with what hazardous-substance surveillance programmes in the United States have long reported, that the substances most often released are common industrial acids and ammonia rather than exotic ones \cite{horton2004,melnikova2020,eiffert2017,hu2004}, and with the older finding that newspaper attention to causes of death is unrelated to their frequency \cite{combs1979}, now measured against an incident registry rather than mortality statistics. We note that our earlier reading of the same six points as a regression of articles per incident on incidents would have suggested a strong inverse relation; that form has a slope of $-1$ under the null \cite{pearson1897}, and the corrected analysis shows no relation at all.

\subsection{Reference bias as a measured quantity}
Selection and description bias \cite{earl2004,strawn2008,zhang2023} concern which events are covered and how. Key-event research \cite{kepplinger1995,brosius1995} shows that a prototypical event raises subsequent attention to similar events. Media templates \cite{kitzinger2000} name the pattern by which a past event becomes the frame for reporting later ones. Reference bias, as measured here, is the quantity that these accounts describe but do not compute: the share of a hazard's coverage that invokes the key event in articles reporting no event at all, measured against a same-block control.\footnote{The term is also used in survey methodology for the bias that arises when respondents rate themselves against different reference groups; the two uses share only the word.} The claim is narrow, and it is worth saying what it is not. It is not a new mechanism. Two established accounts already cover the ground, and neither is displaced here. Availability cascades \cite{kuran1999} describe how an event that is easy to recall becomes, through repetition, the basis of public and regulatory judgement; $\rho$ is one observable cross-section of such a cascade, measured for a single hazard in a single window, and it says nothing about the cascade's dynamics or its regulatory end. The social amplification framework \cite{kasperson1988,pidgeon2003,kasperson2022} describes ripple effects by which an event's signal value outlasts and outgrows its physical consequences, which is close to what a Novichok article citing Salisbury four years later is doing; the difference is that the framework's ripples are traced forward from an event through impacts, whereas $\rho$ is measured backward from present coverage and needs no impact data. Fellenor et al. \cite{fellenor2019} come closest in method, quantifying media volume for a dispersed risk event within that framework; they ask how coverage of one hazard rises and falls over time, where we ask what fraction of a hazard's coverage at any time is about something that already happened, and compare that fraction against a matched control hazard. What is new here is therefore an estimator and a control design, not a mechanism: a number a planner can compute for a hazard on a Tuesday afternoon, where the existing accounts supply the explanation for why the number is not zero. A reader who holds that reference bias is availability cascade under another name loses nothing by saying so, provided the quantity is still computed. A main-subject label misses it, because the past event is usually inserted into coverage of something else; topic taxonomies miss it too, since the taxonomy assigns the article to its hazard regardless of whether the hazard occurred. Measuring it requires an anchor and a control. The same-block control design used here, scoring acids and chlorine against Gumi, is easy to replicate for any hazard with a canonical disaster, and it separates the reference load of the hazard from the general tendency of an outlet or period to look backward.

The Novichok case shows both the size of the effect and the care its measurement needs. Nine articles in ten cite a past poisoning, but the past poisoning most cited is Navalny's, not Salisbury's, and in 2024 the citing articles were reporting his death. The anchor changed over the window as the story acquired a new chapter, which is what collective memory in journalism looks like in practice \cite{edy1999}. The cross-tabulation with the EVENT label is therefore part of the metric, not a check on it: for Novichok, 72\% of anchor mentions occur in articles that report no event, and 65\% of all Novichok coverage cites a past poisoning without reporting one; under the strict reading of Section~\ref{sec:results}, in which the death in custody is not a Novichok event, the figure is 89\%. The size of the effect matters for warning practice. Repeated alerting for hazards that do not materialise erodes public response \cite{sutton2025}, and crisis-alert systems that draw on open sources inherit whatever bias those sources carry \cite{brynielsson2017}; reference bias is the editorial counterpart, the availability heuristic \cite{tversky1973} operating in the newsroom rather than in the reader.

\subsection{Where this sits beside event-based surveillance}
Event-based surveillance is the closest operational practice, and the comparison is worth making precisely because the two are not doing the same thing. An EIOS-family system reads a stream of articles, triages them, verifies a candidate against an independent source and only then records an event; its evaluations report how much of the stream survives that funnel, and for the African region Williams et al.\ \cite{williams2025eios} report both the yield and the delay. Korea runs such a system, and its national evaluation \cite{kim2025ebs} reports the same kind of quantity for communicable disease. What we compute is the yield of the first step alone, for a hazard class where the later steps have nothing to verify against, and across a hazard's whole coverage rather than a live stream. That makes $p_E$ neither a competitor to those figures nor directly comparable with them: a verification pipeline's yield is conditioned on triage having already discarded most of the stream, whereas $p_E$ is measured before any triage. The practical relation is the other way round. Where a verification pipeline exists, $p_E$ says what its input looks like and therefore what its triage step has to remove; where none exists, which is the CBRN case in our setting, $p_E$ is the only available correction between a count and an event estimate, and it is also the reason a pipeline is needed. An agency that already verifies its signals will find little new in the recommendation not to count articles. It may still find the decomposition useful for sizing the hazard classes its pipeline cannot cover.

\subsection{Why topic taxonomies do not help}
Adding the archive's incident tags to the classifier did not improve it. This is the same finding as the weak convergence between the tags and the EVENT label ($\phi = 0.11$), seen from the other side. The taxonomy is good at its job, since ``nuclear accident'' articles are indeed about nuclear accidents, but 90\% of them report no accident. Any surveillance system that inherits a news taxonomy as its event definition inherits this conflation. The character $n$-gram model, by contrast, learns the lexical markers of occurrence (times, places, quantities, casualty counts, agency responses) from 200 characters of text, and does so with a footprint small enough to run continuously on modest hardware, and it transfers to accident and terrorism queries it was not trained on with a loss of a few points.

The classifier performs worst in the security block, which is where the decomposition matters most. Security-block EVENT articles are few and lexically diverse; the model's EVENT-$F_1$ there is 0.51 in cross-validation and 0.40 on unseen queries, and it misses more than half of the events. Two remedies are available without changing the design: open-set alerting, which flags articles the model has not seen the like of rather than forcing them into a class \cite{kang2026openset}, and the use of structured incident records as a supervision signal where they exist, as has been done for asymmetric threat prediction from open-source incident logs \cite{bae2025balloon}. For the security block, post-stratification of coded samples is in any case the better route to a corpus estimate, since it needs no sensitivity correction; the classifier's role there is to triage, not to count.

\subsection{Limitations}
\label{sec:limits}
Coding was performed by two language models from the same vendor; their agreement was substantial and highest on the boundary that matters most, but agreement between related models bounds reliability from above and says nothing about validity. The human validation of Section~\ref{sec:human} shows what that distinction costs: the models agreed with each other at $\kappa = 0.79$ while agreeing with the human standard at 0.50 on the EVENT boundary, and at 0.24 in the security block, where they counted topically adjacent occurrences as events. Model labels are used throughout the results because they cover all 2,867 articles, but every claim about the size of the security event share rests on the 400-article human subset, and the corpus correction is only as good as the sensitivity and false-positive rate estimated from it, which for the security block carry intervals of 0.28--0.83 and 0.08--0.10. The subset was coded by two of the authors, who were not blind to the study's hypothesis, although they were blind to the model labels and to each other, worked from a codebook fixed before the subset was drawn, and saw the articles in shuffled order; no third coder adjudicated the 40 disagreements, which is why the estimate is reported across treatments of them rather than as a single number. The human coding form recorded the component but not the event date, place and object. For the security block one of us supplied those fields for the 22 human-confirmed events in a third pass, by a single analyst who already knew the results; it is reported with its clusters deposited rather than treated as independent coding, and the reported amplification keeps the conservative model-based duplication factor. A replication should collect the fields at coding time. The anchor lexicons were written by the authors, and $\rho$ depends on the text basis: the keyword field roughly doubles every value relative to title and excerpt, which is why both are reported and only the contrast with controls is interpreted. The study covers one country, one language and one archive over a window that contains the war in Ukraine, which dominates the ``drone attack'' query and part of the conventional block. The duplication factor $m$ is bounded by the sampling design in the gold set and by an approximate text rule in the corpus, and the text rule was tuned against the clusters it is then compared with, so its agreement is an upper bound; the block contrast rests on $p_E$, which depends on neither. The registry comparison rests on six substances, too few to identify a severity-adjusted model from the casualty and quantity fields the registry also holds. Anchors were specified for eight queries and without a within-block control for the security block. The DPRK nuclear-test query dominates the security corpus and is geopolitical coverage; the gold set weights it as one query in thirty and the contrasts survive its exclusion, but corpus-level security estimates should be read with it in mind. Finally, the archive supplies only the first 200 characters of body text, which bounds what any coder or model can see and biases the title-and-excerpt $\rho$ downward, since a mention of the anchor later in the article is missed.

\subsection{A procedure for practice}
\label{sec:practice}
The Sendai Framework's first priority is understanding disaster risk, and it asks national authorities to build and maintain hazard and loss data on which everything downstream rests; its Target G asks them to extend multi-hazard early warning. Both run into the same wall for the hazards considered here. There is no loss database for a nerve-agent attack that did not happen, and an early-warning system needs a baseline rate to alert against. Open sources are what agencies reach for when the registry is missing, and the progress reviews of multi-hazard early warning note that the data underpinning such systems is often the weakest link rather than the technology \cite{rokhideh2025mhews}. The decomposition below is a way of stating what an open-source baseline is actually made of before it is used as one. It does not create incidence data where none exists; it says how much of a volume figure is not incidence, which is the prior question for both Priority 1 and Target G.

The steps below apply these results to a single hazard and a single archive. The cost is one coding pass over a few hundred headlines, which a single analyst can do in a day or two per hazard; running it as we ran it, with two independent coders on 400 articles, doubles that and is what we would recommend for a hazard whose number will be quoted. A lexicon of twenty hazards is therefore a few analyst-weeks, not an afternoon, and it is a one-off per hazard rather than a recurring load, since the quantities are properties of a hazard's coverage over a window rather than of a live stream. Where that cost is not available, coding 200 articles with one coder gives the event share to within about seven points at the values we observe, which is enough to separate a hazard whose coverage is mostly events from one whose coverage is mostly not.

\emph{Fix the query and the window before looking at any counts.} Write the query from an external list rather than from what returns articles: a registry's substance list, a statutory hazard schedule, an incident taxonomy. Record it verbatim. A query written to maximise yield will maximise the off-topic share, which is the component that later inflates the count.

\emph{Draw a sample and code it.} Two hundred articles per hazard is the working minimum and 400 is comfortable. Sample within strata if the archive allows it, and take every article a first-pass filter calls an event, since that is the stratum in which errors are found. Code each article into the five components of Eq.~\eqref{eq:identity} from title and lead alone, applying the decision order and the means-mismatch rule; the codebook released with this article is usable as it stands. Use two coders where possible and report their disagreement rate: give the estimate under each coder separately and under the extreme resolutions of their disagreements, as in Section~\ref{sec:human}, and treat the width of that range rather than a bootstrap interval as what the figure is known to within.

\emph{Report four quantities.} The event share $p_E$ with its interval; the same share among on-topic articles only, which separates the hazard from the query; the domestic event share, since a hazard's foreign coverage does not describe the risk at home; and the amplification factor $A$, which converts the raw count into a bounded event estimate. For an event share near 10\%, 100 coded articles give a 95\% interval of about 6--17\%, 200 give 7--15\% and 400 give 7--13\%; near 30\%, the same samples give 22--40\%, 24--37\% and 26--35\%.

\emph{Benchmark where a registry exists, and measure reference bias where one does not.} A registry comparison bounds what the coverage can be read as, though it establishes only that the two diverge. Where no registry exists, write an anchor lexicon for the hazard's canonical disaster and one for a neighbouring hazard in the same sector, and compare the citation rates on the same window. A hazard whose coverage cites a past disaster at several times the rate of its neighbours is carrying inherited salience, and the share of that coverage reporting nothing new is the quantity to put in front of a decision.

\emph{Stop when the query or the archive cannot answer the question.} Two thresholds are worth stating, and one of them does not do the work one might expect. If the on-topic share is below about a third, the query is retrieving the wrong thing and should be rewritten before anything is estimated; this catches a badly written query but it does not catch the case this paper is about. Our own security block passes it comfortably, at 74\% on topic, and the coverage is still not an incidence signal, because what the query retrieves is genuinely about the hazard and still reports no event. If the event share is low and its coverage is largely foreign, the count is describing the archive's reporting threshold as much as the hazard, and no correction recovers what was never reported; the output should then be a bound, with a statement that the news cannot answer the question.

The identity in Eq.~\eqref{eq:identity} and the procedure above are not specific to Korea, to CBRN hazards or to this archive. The numbers are. Coverage volume responds to the observing country's press density, wealth and political interests as much as to what happened \cite{galtung1965,vanbelle2000,zhu2018}, and Korea is at the top of those distributions with 54 national outlets in a single searchable archive. In a market with fewer outlets, the duplication factor $m$ will be smaller and the amplification correspondingly milder; where a hazard's events are mostly domestic rather than foreign, the event share should be higher than the 5\% reported here, since the reporting threshold that suppresses ordinary foreign occurrences does not apply. What we would expect to hold is the ordering: a hazard that is rare, feared and institutionally active will have a lower event share and a higher reference load than a common hazard of the same sector, because both quantities follow from the same scarcity of new events and abundance of institutional and retrospective material. That expectation is a hypothesis for replication, not a result of this study.

%% ------------------------------------------------------------------
\section{Conclusion}
\label{sec:conclusion}
News volume for hazard queries is mostly not event coverage. Across the gold set, which weights each of the 30 queries equally, one article in three reported an event, and post-stratified onto the 109,610-article corpus the figure is one in four; for CBRN security hazards, one in ten by model coding and one in twenty (5.0\%, 95\% CI 2.5--7.9) against a two-coder human reference standard built on the 219 articles of this block that the two coders agreed on, and, on the same human standard, about one in forty reported a domestic event, a figure that is stable only because the model's errors in that count run in both directions. The event share separates CBRN hazards from both control blocks, which could not be distinguished from each other at ten queries per block, and the contrast survives the exclusion of every query a reader might assign differently. On the human labels the separation holds at block level with query-clustered intervals, and the omnibus rank test still rejects; the per-query rank comparison with the industrial block does not, which is a limit of a subset holding a median of seven agreed articles per control query rather than a reversal. Reference bias, defined as citation of a past anchor event, reaches 33--89\% for hazards with such an event and 1--2\% for matched controls, most of those citations occur in articles that report nothing new, and a rule that names no anchor reproduces the gap, so the control values record the absence of a canonical occurrence rather than the absence of the anchor we scored. Against the official registry, article counts for six chemicals show no detectable relation to incident counts, and the framework's event estimates exceed the registry's for every substance. A lightweight classifier recovers the event component for accident and terrorism queries, including unseen ones, but not for CBRN queries, and the archive's own taxonomy does not help. Each of these results is a reason not to count articles. Together they supply the arithmetic for converting a volume count into an event estimate, which is what the raw count was taken to be.

%% ------------------------------------------------------------------
\section*{CRediT authorship contribution statement}
\textbf{Miae Kim:} Investigation, Data curation, Validation, Writing -- review \& editing. \textbf{Janghyeon Cha:} Resources, Investigation, Validation, Writing -- review \& editing. \textbf{Myeongsik Shin:} Software, Formal analysis, Visualization, Writing -- review \& editing. \textbf{Gyungjun Choi:} Investigation, Data curation, Writing -- review \& editing. \textbf{Ku Kang:} Conceptualization, Methodology, Software, Formal analysis, Validation, Writing -- original draft, Supervision, Project administration.

\section*{Declaration of competing interest}
The authors declare no competing financial or personal interests.

\section*{Declaration of generative AI and AI-assisted technologies}
Two large language models (Claude Opus and Claude Sonnet, Anthropic) were used as independent coders under a written codebook, as described in Section~\ref{sec:coding}. Their agreement with each other and with two human coders is reported in Sections~\ref{sec:clf} and~\ref{sec:human}, and the manuscript states which results the model labels support and which they do not. Language models were also used to assist with drafting; the authors reviewed and edited all content and take full responsibility for it.

\section*{Funding}
This research did not receive any specific grant from funding agencies in the public, commercial, or not-for-profit sectors.

\section*{Acknowledgements}
The authors thank the Korea Press Foundation for access to the BigKinds archive and the National Institute of Chemical Safety for the public chemical-accident registry.

\section*{Data availability}
The article corpus was retrieved from BigKinds (Korea Press Foundation) under its terms of use and cannot be redistributed in full. Article identifiers for all 109,610 articles, the gold-set labels with event identifiers, the labels of both human coders and the verification-sample design, the codebook, the anchor dictionary, the chemical-accident registry extract, the seed values and all analysis and figure code are deposited at \url{https://github.com/bisu9082/refbias}, released under the MIT Licence for code and CC BY 4.0 for labels and derived outputs. A reader with BigKinds access can rejoin the article text on the deposited identifiers; every result computed from labels rather than text is reproducible without that access, and the deposit includes the three verification scripts used to check the manuscript's numbers against the deposited outputs: a regression check on pinned values, a numerical check of the published tables, and a three-layer audit that recomputes the headline quantities from the raw labels, matches them against the manuscript text and cross-checks them between output files. The wrapper text that presented each article to the model coders, and the model identifiers returned by the interface, were not retained. The registry is available from the Korean open-data portal (dataset 15069200). The analysis was re-run end to end on an independent Ubuntu 24.04 installation with different versions of NumPy (2.5 against 2.4), SciPy (1.18 against 1.17) and scikit-learn (1.9 against 1.8), pandas being 3.0 in both, and reproduced every value reported here to the precision at which it is given. Results that depend on the article text rather than on the deposited labels, which are the anchor rates of Section~\ref{sec:refbias} and the classifier, cannot be recomputed by a reader without BigKinds access; the query, window and fields needed to retrieve the same records are given in Section~\ref{sec:methods}.

\bibliographystyle{elsarticle-num-names}
\bibliography{myref}

\end{document}
